% P03.1 finite-orbit symmetry defect proof, live integration version.

\begin{proof}[\Cref{thm:p031-invariant-real-pl-descent}]
Pullback sends every member of $\mathcal P_X(f)$ to a $G$-invariant member
of $\mathcal P_E(f\circ\pi)$, and it is injective because $\pi$ is
surjective.

Conversely, let $v\in\mathcal P_E(f\circ\pi)^G$.  Every fiber of $\pi$ is a
single $G$-orbit, so $v$ is constant on fibers and there is a unique
set-theoretic function $b:X\to\mathbf R$ such that
\begin{equation}
\label{eq:p031-invariant-real-pl-descends}
                              v=b\circ\pi.
\end{equation}
Apply \cite[Proposition~5.11]{GJR26} with the divisible coefficient group
$\mathbf R$.  Since $\pi$ is finite, flat, and surjective, the proposition
shows directly that $b$ is real PL.  Semipositivity reflects through the
proper surjection $\pi$ by
\cite[Proposition~7.6\textup{(1')}]{GJR26}.  Finally,
$v\leq f\circ\pi$ and surjectivity give $b\leq f$.  Thus
$b\in\mathcal P_X(f)$, and the construction is inverse to pullback and
preserves pointwise order.

Now let $u\in\mathcal P_E(f\circ\pi)$.  Every translate $u\circ g$ is
semipositive real PL.  Such functions form a convex cone by
\cite[\S7.3\textup{(5)}]{GJR26}, so $\mathsf A_Gu$ is semipositive real PL.
It is $G$-invariant and remains below $f\circ\pi$.  This proves the averaging
claim.  A nonempty cover family therefore has a nonempty invariant subfamily
and hence a nonempty base family; the converse follows by pullback.  Taking
pointwise suprema in \eqref{eq:p031-invariant-real-pl-bijection} proves
\eqref{eq:p031-invariant-real-pl-subenvelope}, including the simultaneous
empty-family case.
\end{proof}

\begin{proof}[\Cref{thm:p031-centered-symmetry-defect}]
A base competitor and its pullback bound both raw envelopes from below.
They also force $f$ to be real-valued, so the obstacle bounds both envelopes
from above.  Hence the two envelope values and their difference are finite.
Pullback of base competitors gives
\[
                 s_{\mathrm{PL},X}\circ\pi
                     \leq s_{\mathrm{PL},E},
\]
and therefore $\Sigma_{\pi,f}\geq0$.

For $g_0\in G$, composition $u\mapsto u\circ g_0$ permutes the cover
competitor family.  Consequently
\[
 s_{\mathrm{PL},E}(g_0(e))
  =\sup_{u\in\mathcal P_E(f\circ\pi)}u(g_0(e))
  =s_{\mathrm{PL},E}(e).
\]
The pullback term in \eqref{eq:p031-pl-symmetry-defect} is invariant as
well.  Thus the defect is invariant and, since each fiber is one orbit,
constant on fibers.

By \Cref{thm:p031-invariant-real-pl-descent}, the average
$\mathsf A_Gu$ is an invariant cover competitor.  Equation
\eqref{eq:p031-invariant-real-pl-subenvelope} therefore gives
$\mathsf A_Gu\leq s_{\mathrm{PL},X}\circ\pi$, so $H_u\geq0$.  For every
cover competitor, there is the exact pointwise cancellation
\[
 u(e)-\mathsf A_Gu(e)-H_u(e)
   =u(e)-s_{\mathrm{PL},X}(\pi(e)).
\]
Taking the supremum proves
\eqref{eq:p031-centered-symmetry-formula}.

Discarding the nonnegative headroom gives
$\Sigma_{\pi,f}\leq\Omega_{\pi,f}$.  The pullback of any base competitor is
invariant and has $u-\mathsf A_Gu=0$, so $\Omega_{\pi,f}\geq0$.  Moreover,
for each individual competitor,
\[
 u(e)-\mathsf A_Gu(e)
   \leq
      \max_{g\in G}u(g(e))-\min_{g\in G}u(g(e)).
\]
Taking suprema proves \eqref{eq:p031-symmetry-oscillation-bounds}.  The
centered formula is a supremum of its displayed summands, while the defect
is nonnegative.  It vanishes exactly when every summand is nonpositive,
which is \eqref{eq:p031-pointwise-symmetry-criterion}.

Finally, every locally Fubini--Study function is semipositive real PL by
\cite[Definition~7.8 and \S7.9]{GJR26}.  Adjoining LFS minorants to the
semipositive real-PL family therefore does not enlarge it.  The P03 union
envelopes are precisely the two PL envelopes, and
\eqref{eq:p031-full-union-symmetry-defect} follows from the definition of
$\Sigma_{\pi,f}$.  Its vanishing criterion is immediate, and the sufficient
condition $\Omega_{\pi,f}\equiv0$ follows from
\eqref{eq:p031-symmetry-oscillation-bounds}.  None of these pointwise
identities supplies any of the regularity or vanishing statements excluded
in the theorem.
\end{proof}
